% main.tex -- primary category gr-qc, cross-list astro-ph.HE
% "Geodesic: a parallel test-particle geodesic tracker for the Einstein Toolkit"
% Compile: pdflatex main && bibtex main && pdflatex main && pdflatex main
\documentclass[11pt]{article}

\usepackage[T1]{fontenc}
\usepackage[utf8]{inputenc}
\usepackage{amsmath}
\usepackage{amssymb}
\usepackage{graphicx}
\usepackage{booktabs}
\usepackage{natbib}
\usepackage{float}
\usepackage[colorlinks=true,linkcolor=blue,citecolor=blue,urlcolor=blue]{hyperref}

\setlength{\textwidth}{16.5cm}
\setlength{\textheight}{22.5cm}
\setlength{\oddsidemargin}{-0.2cm}
\setlength{\evensidemargin}{-0.2cm}
\setlength{\topmargin}{-1.2cm}

% figures: hard error if a file is missing (no silent placeholder)
\newcommand{\maybefig}[2]{\includegraphics[width=#2]{#1}}

% data values (single source of truth; filled after the verification runs)
\newcommand{\num}[1]{\csname gnum#1\endcsname}
\newcommand{\gnumconsexp}{9}
% math values: used only inside $...$ at the call sites
\newcommand{\gnummaxdE}{2.8\times10^{-10}}
\newcommand{\gnummaxnorm}{8.9\times10^{-16}}
\newcommand{\gnumexacterr}{1.7\times10^{-6}}
\newcommand{\gnumconvorder}{5.6}
\newcommand{\gnumscaleeff}{10.4}
\newcommand{\gnumhydrofrac}{3.1}

\title{Geodesic: a parallel test-particle geodesic tracker for the Einstein Toolkit}

\author{Yuhua Li\thanks{liyhts@yeah.net} \\
% The author is retired, so no institutional affiliation is given
{\small Independent researcher}}

\date{\today}

\begin{document}

\maketitle

\begin{abstract}
We present \texttt{Geodesic}, a new thorn for the Einstein Toolkit that
integrates test particles on geodesics of a spacetime in Kerr--Schild
Cartesian coordinates.  The thorn provides two modes of metric
evaluation: an exact analytic Kerr--Schild metric, independent of the
computational grid and verified symbolically to be a vacuum solution,
and a grid mode that reads the metric from the simulation through one
of four local samplers selected by a single parameter (\texttt{l27},
\texttt{local27}, \texttt{aei}, \texttt{carpet}).  Each particle is an
independent unit of work, so the particle loop scales to $10^4$
particles, reaching a measured $9.2\times$ speedup on the 88 logical
threads of a 44-core node; the same independence underlies an
owner-computes MPI mode with the replicated state bit-identical on
every rank, and the samplers run unchanged on Carpet AMR meshes.  The
geodesic equation is integrated with an adaptive
Runge--Kutta--Fehlberg (RKF45) scheme and a per-step renormalization of
the four-velocity.  We verify the thorn by (i) conservation of the
energy $E$, the axial angular momentum $L_z$ and the norm to
$\sim 10^{-\num{consexp}}$ over $T=1000\,M$ on three bound orbits;
(ii) grid convergence of the default sampler with a measured local
order of \num{convorder}; (iii) a comparison of the grid samplers on a
precessing orbit of $100$ particles; (iv) stable orbits on a 3-level
AMR mesh; and (v) robustness on the two-puncture \texttt{GW150914}
background: geodesics integrate stably on the \texttt{ML\_BSSN}-evolved
(time-dependent) merger spacetime through $t = 14.05\,M$, with the stability
thresholds and costs of background resolution and timestep quantified; on the
frozen initial data the strong-field drops of a near-ring orbit are
deterministic and essentially back-end-independent, and a three-resolution scan
gives the observed-order trend of each sampler.  Embedded in a full \texttt{GRHydro}
evolution, the integration of $10^4$ particles adds only a few percent
to the total wall time.  The thorn is released under the GPLv2+ license
together with the configuration files used in this work.
\end{abstract}

\section{Introduction}

Test-particle geodesics are a basic diagnostic in numerical relativity and
in modelling of accretion flows around black holes \citep{poisson2004}:
they provide reference
orbits against which self-consistent hydrodynamic and electromagnetic
solutions can be compared \citep{bz1977}, they define the orbital structure used to
interpret line profiles and disk precession \citep{bardeen1973,gralla2020}, and they serve as a
validation tool for the spacetime produced by a simulation
\citep{kerr1963,bardeen1972,wilkins1972}.
General-purpose geodesic integrators for 3+1-split spacetimes have been
developed for test-particle diagnostics in GRMHD simulations
\citep{bacchini2018,bacchini2019}; dedicated packages compute light
propagation and optical observables on numerical-relativity output
\citep{grasso2022} and polarized ray tracing on prescribed metrics
\citep{aimar2024}; timelike and plunging Kerr geodesics are available
analytically or through public packages \citep{park2024,dyson2023}; and in
self-force work the particle worldline is advanced together with the field
on the numerical grid \citep{vretinaris2026}.  None of these, however, is a
test-particle geodesic tracker inside the Cactus schedule: the verification
studies of \citet{bacchini2018,bacchini2019} target fixed or given metrics,
and none of these codes addresses particle tracking on an evolved spacetime
resolved on an AMR mesh distributed across MPI ranks.  The standard
frameworks for
general-relativistic simulations of compact objects, the Cactus code
framework \citep{cactusweb} and its extension, the Einstein Toolkit
\citep{loeffler2012}, do not ship a geodesic tracker for test particles in
Kerr--Schild coordinates, which is the coordinate system used by most
Kerr initial-data and spacetime-evolution thorns.  Users who need geodesics
today must either write their own integrator or post-process the evolved
metric fields, and in both cases the cost of a global interpolation of the
metric at every particle position prevents the straightforward parallel
scaling that the problem structure allows.

\texttt{Geodesic} is a thorn for the Einstein Toolkit that supplies this
functionality.
Its design goals are (i) to support both an exact analytic Kerr--Schild
metric and a metric interpolated from the simulation grid, so that the same
particle set can be tracked on the analytic background or on an evolved
spacetime; (ii) to keep the per-particle work completely local, so that the
particle loop can be parallelized with OpenMP without any inter-particle
communication; and (iii) to be a clean, auditable reference implementation
of the geodesic equation in Kerr--Schild Cartesian coordinates, verified
against an independent implementation.

The second goal is a statement about data dependencies, not about achieved
efficiency: the loop has no inter-particle communication, but its measured
OpenMP scaling is limited by memory and cross-socket bandwidth rather than
by load balance, and we report that number rather than an idealised one
(Sec.~\ref{sec:scaling}).  A parallelization that co-locates particles with
the grid data they read would be expected to do better; we have not
implemented one.

The main contributions of this paper are:
\begin{enumerate}
\item a thorn that integrates the geodesic equation in Kerr--Schild
  Cartesian coordinates with an adaptive RKF45 scheme and a per-step
  four-velocity renormalization;
\item an exact analytic Kerr--Schild metric whose vacuum character is
  verified symbolically (Sec.~\ref{sec:exactmetric});
\item four local grid-metric samplers selected by one parameter
  (\texttt{l27}, \texttt{local27}, \texttt{aei}, \texttt{carpet};
  Sec.~\ref{sec:samplers}), in which every particle is an independent
  unit of work, validated by a strong-scaling run with $10^4$ particles
  (Sec.~\ref{sec:scaling});
\item a verification suite: conservation of $E$, $L_z$ and the norm over
  $T=1000\,M$ for three distinct bound orbits, a grid-convergence study
  showing a measured local convergence order of \num{convorder}
  for the default sampler, a four-way comparison
  of the exact metric against the grid samplers on a precessing orbit of
  $100$ particles, and a robustness run on the two-puncture GW150914
  merger background (Sec.~\ref{sec:gw150914}).
\end{enumerate}

The remainder of the paper is organized as follows: Sec.~\ref{sec:method}
describes the geodesic equation in the 3+1 formulation and the
initialization of the particle four-velocity; Sec.~\ref{sec:metric} the
exact mode and the four grid-metric samplers; Sec.~\ref{sec:impl} the
implementation, including
the OpenMP and MPI parallelization, the AMR level handling and the
particle-lifecycle management;
Sec.~\ref{sec:verification} the verification results; Sec.~\ref{sec:scaling}
the strong-scaling results; and Sec.~\ref{sec:concl} concludes.

\section{Method}\label{sec:method}

\subsection{Geodesic equation}

We work in Kerr--Schild Cartesian coordinates $(t,x,y,z)$ with spin axis
along $+z$, in which the Kerr metric of mass $M$ and dimensionless spin
$a$ takes the form \citep{kerrschild1965}
\begin{equation}
g_{\mu\nu} = \eta_{\mu\nu} + \frac{2Mr^3}{D}\,l_\mu l_\nu ,
\qquad
l_\mu = \left(1,\, \frac{a y + r x}{Q},\, \frac{r y - a x}{Q},\, \frac{z}{r}\right),
\label{eq:metric-form}
\end{equation}
with $\eta_{\mu\nu}=\mathrm{diag}(-1,1,1,1)$,
$Q = r^2 + a^2$, $D = r^4 + a^2z^2$, and $r$ the Kerr radial coordinate
defined by $r^2 = (x^2+y^2+z^2-a^2 + \sqrt{(x^2+y^2+z^2-a^2)^2 + 4a^2z^2})/2$.
The vector $l_\mu$ is null, $g_{\mu\nu}l^\mu l^\nu = 0$, and the metric is a
vacuum solution, $R_{\mu\nu} = 0$ (Sec.~\ref{sec:exactmetric}).

The four-velocity $u^\mu = \mathrm{d}x^\mu/\mathrm{d}\tau$ of a massive
test particle obeys the geodesic equation
\begin{equation}
\frac{\mathrm{d}u^\mu}{\mathrm{d}\tau} =
-\Gamma^\mu_{\alpha\beta}\,u^\alpha u^\beta ,
\qquad
g_{\mu\nu}u^\mu u^\nu = -1 ,
\label{eq:geo}
\end{equation}
where $\Gamma^\mu_{\alpha\beta}$ are the Christoffel symbols of $g_{\mu\nu}$.
Because the metric is stationary and axisymmetric, two quantities are
conserved along the geodesic: the specific energy
$E = -u_t$ (Killing vector $\partial_t$) and the axial angular momentum
$L_z = -y\,u_x + x\,u_y$ (Killing vector $-y\partial_x + x\partial_y$),
with the covariant components $u_\mu = g_{\mu\nu}u^\nu$.  Both
are used as verification monitors.

\subsection{Integration in proper time}

The thorn integrates Eq.~(\ref{eq:geo}) in proper time $\tau$ with an
adaptive fifth-order Runge--Kutta--Fehlberg scheme (RKF45,
\citealt{fehlberg1969}; the \texttt{rkf45} driver of the GNU Scientific
Library~\citep{gsl}), with an absolute error tolerance of $10^{-8}$ and
zero relative tolerance.  Each Cactus time step advances the particle until
its coordinate time reaches $t + \Delta t$: the proper-time increment is
predicted so that $t$ lands on the target, and any overshoot is removed by
a linear back-interpolation, so that all particles share the Cactus time
axis.  The number of proper-time sub-steps within one Cactus step is
bounded by a per-particle cap (\texttt{max\_iter\_cap}, default 2000), so
that a particle whose proper-time step has collapsed (as near a horizon)
cannot stall the evolution.  At the end of every Cactus time step the
thorn renormalizes the
four-velocity so that $g_{\mu\nu}u^\mu u^\nu = -1$ holds to roundoff,
which removes the secular drift of the norm that pure adaptive integrators
accumulate.  The spatial components $u^i$ produced by the ODE step are kept
as they are, and the constraint is solved for $u^t$ as a quadratic,
\begin{equation}
g_{00}\,(u^t)^2 + 2\,g_{0i}\,u^i\,u^t + \left(1 + \gamma_{ij}u^i u^j\right) = 0,
\label{eq:renorm}
\end{equation}
selecting the future-directed root.  Outside the ergosphere
($g_{00}<0$) the two roots have opposite signs and the unique positive
root is unambiguous; inside the ergosphere the same $u^i$ can
in fact correspond to two future-directed solutions (a geometric effect
of frame dragging), and the root is chosen by continuity, i.e.\ the one
closest to the $u^t$ produced by the ODE step before renormalization.

\subsection{Initialization of the four-velocity}

The particle initial conditions are specified in the par file as a position
$(x,y,z)$ and the contravariant four-velocity components
$u^\mu=(u^t,u^x,u^y,u^z)$ in Kerr--Schild coordinates.  At the first time
step the thorn evaluates the metric at the initial position and checks the
normalization $g_{\mu\nu}u^\mu u^\nu = -1$.  If it holds to a tolerance of
$10^{-6}$, the input is used unchanged and the log says so; otherwise $u^t$
is re-solved from Eq.~(\ref{eq:renorm}) with $u^i$ held fixed, and the
correction (before/after values) is written to the log.  The check is
therefore idempotent: feeding the values logged at initialization back into
the par file changes nothing.  The $10^{-6}$ tolerance is loose on purpose;
the per-step renormalization described above restores the constraint to
roundoff on the first step regardless.  All orbits used in this paper are
generated by first integrating the exact-metric geodesic equation with an
independent Python implementation (RK4, $\Delta t = 0.02\,M$) and then
writing the resulting exact four-velocity directly into the par file; the
initial data are listed in Table~\ref{tab:orbits}.

\section{Metric evaluation}\label{sec:metric}

\subsection{Exact Kerr--Schild mode}\label{sec:exactmetric}

With \texttt{Exact = yes} the thorn evaluates the Kerr--Schild metric and
its coordinate derivatives analytically at the particle position
(Eq.~(\ref{eq:metric-form})), so that the result is independent of the computational grid.
The analytic Christoffel symbols are obtained from the analytic first
derivatives of the metric; no finite differences enter the evolution.  As a
self-consistency check we verify the vacuum character of the implemented
metric symbolically with SymPy~\citep{sympy}: the components of the
thorn's metric are checked term-by-term against the Kerr--Schild form
Eq.~(\ref{eq:metric-form}) (which also proves $l_\mu$ is null), and the
Ricci tensor $R_{\mu\nu}$ is assembled from the analytic metric and its
analytic inverse $g^{\mu\nu} = \eta^{\mu\nu} - H\,l^\mu l^\nu$ and evaluated
at five generic points to 40-digit precision, giving
$\max|R_{\mu\nu}| < 10^{-170}$, i.e.\ the implemented metric is a vacuum
solution to beyond machine precision.

\subsection{Grid-metric samplers}\label{sec:samplers}

With \texttt{Exact = no} the metric is taken from the simulation grid: the
thorn reads the \texttt{ADMBase} lapse, shift and spatial metric (declared
with schedule \texttt{READS} clauses), assembles the covariant metric
$g_{\mu\nu}$ at the stencil grid points surrounding the particle, and
obtains the spatial metric derivatives there by central finite
differences (sixth order in the stencil-based samplers).  The
Christoffel symbols also require the time derivative
$\partial_t g_{\mu\nu}$; how it --- and the metric value at the particle's
coordinate time --- is obtained from the stored \texttt{ADMBase} time
levels is fixed per sampler below.  A stationary metric is detected
automatically (the oldest two stored levels agree to relative $10^{-12}$
on a subsample), in which case $\partial_t g_{\mu\nu} = 0$ exactly and the
time differencing is skipped.  Within one particle step the connection is
frozen at the grid's current time, so for a time-dependent metric the
particle step is capped at the grid time step; time-dependent runs
therefore need the \texttt{ADMBase} fields at three time levels
(\texttt{Carpet::init\_fill\_timelevels = yes}).  How the metric and the
Christoffel symbols are sampled at the (generally off-grid) particle
position is selected by the parameter \texttt{metric\_interp}; the four
choices are described below.  All four samplers are local: the
computation for one particle reads only the grid cells around it and
writes no data used by any other particle.

\paragraph{l27.}
The simplest sampler: for a given particle the thorn identifies the 27
grid points of the 3$\times$3$\times$3 cell stencil surrounding it and,
for each metric component and each Christoffel symbol, fits the unique
polynomial that is at most quadratic in each coordinate and interpolates
the 27 stencil values --- a 27-term tensor-product Lagrange interpolant
\citep{berrut2004}.  The fit is a single matrix--vector product with a
precomputed 27$\times$27 inverse Vandermonde matrix, and fields that are
constant on the stencil collapse to their constant value without the
product.  The nodal Christoffel symbols are computed from the sixth-order
finite-difference metric derivatives, and $\partial_t g_{\mu\nu}$ from a
second-order backward difference $(3L_0-4L_1+L_2)/(2\,\Delta t)$ over the
three stored time levels (zero for a stationary metric); the metric values
are taken from the current time level only.  The fit is cached per
particle: it is built once per Cactus time step (and re-evaluated within
the step only if the particle crosses into an adjacent stencil cell), and
every subsequent adaptive RKF45 substep and stage evaluates the cached
polynomials alone --- within a step the right-hand side touches neither
the grid nor the finite differences.  The fit is node-exact on the 27
stencil nodes and, for smooth fields, locally third-order accurate,
$\mathcal{O}(\Delta x^3)$, while the finite-difference error of the
stencil values is $\mathcal{O}(\Delta x^6)$ and does not limit the
convergence.  A degenerate or non-finite fit drops the particle via the
loss policy (Sec.~\ref{sec:impl}); there is no fallback interpolant in
this mode.  The sampling points are clipped to the cells interior to the
local processor box, so the stencil never reaches the ghost region and the
sampler is safe under MPI domain decomposition; in a multi-process run a
particle sitting on, or within a few cells of, a rank face has its stencil
centre clipped away from it, so the one-sided 27-point fit fails its
validity check and the particle is dropped at the first step (the other
back-ends are unaffected) --- with the default $2\times2$ ($y,z$)
decomposition of a 4-rank unigrid run every equatorial particle lies on
the $z=0$ rank face, which is why the testsuite runs \texttt{l27}
single-process.

\paragraph{local27 (default).}
A self-contained variant of \texttt{l27} with a wider stencil and a
higher-degree fit: for each of the three stored time
levels the thorn fits, for each metric component and each nodal
Christoffel symbol, the 125-node ($5\times5\times5$) quartic
tensor-product polynomial through the stencil surrounding the particle
(three separable $5\times5$ inverse-Vandermonde transforms, the
$125\times125$ analogue of the \texttt{l27} matrix), and blends the three
level fits with a three-node (quadratic) Lagrange interpolation in the
particle's coordinate time.  The fit is node-exact on all 125 grid nodes
and formally fifth-order accurate in space.  It uses no external
interpolator thorn.  The
fits are cached per particle as in \texttt{l27}.  Two degradation paths
keep the sampler well defined on AMR meshes and in strong curvature: when
a past time level is not yet filled (the first step, or a freshly grown
refinement region) the time blend degrades to a two-level linear blend or
to the current level only; and when the nodal Christoffel samples are
non-finite (a sixth-order finite-difference window touching the excised
or singular region) the Christoffel symbols are instead assembled
analytically from the metric fits and their spatial derivatives.

\paragraph{aei (fourth-order pointwise).}
\texttt{metric\_interp = aei} evaluates, at the particle position, a
fourth-order pointwise Lagrange interpolant of the ten 3+1 fields (lapse,
three shifts, six spatial-metric components) and of their first spatial
derivatives, using the \texttt{AEILocalInterp} thorn
\citep{aeithornguide}; the metric and its Christoffel symbols are then
composed pointwise from these sampled values.  The two most recent stored
time levels are blended linearly in the particle's coordinate time, and
$\partial_t g_{\mu\nu}$ is the first-order difference
$(L_0-L_1)/\Delta t$.  Unlike \texttt{l27} and \texttt{local27}, no
per-particle fit is cached: the interpolant is re-evaluated at every
RKF45 stage.

\paragraph{carpet.}
\texttt{metric\_interp = carpet} is a pointwise sampler with the same
fourth-order \texttt{AEILocalInterp} spatial molecule as \texttt{aei},
evaluated on the current time level; the time derivative is the
second-order three-point difference $(3L_0-4L_1+L_2)/(2\,\Delta t)$ of the
full 4-metric, the convention of the \texttt{CarpetInterp} driver
\citep{carpethornguide} after which the sampler is named.  Like
\texttt{aei}, it runs unchanged under OpenMP and the owner-computes MPI
mode.

\section{Implementation}\label{sec:impl}

\texttt{Geodesic} is written in C as a Cactus/Einstein Toolkit thorn.
The particle state
$(t,x,y,z,u^t,u^x,u^y,u^z)$ is stored in Cactus interface arrays
\texttt{particle\_tt}, \texttt{particle\_tx}\ldots, and the integration loop
is a plain C loop over the particles, wrapped in an OpenMP
\texttt{parallel for} with \texttt{schedule(dynamic,1)}:
\begin{verbatim}
#pragma omp parallel for schedule(dynamic,1)
for (mi = 0; mi < particle_n_total; mi++) {
    /* per-particle RKF45 step on the locally evaluated metric */
}
\end{verbatim}
There is no shared state inside the loop: each thread evaluates the metric
for its own particle (either analytically, Sec.~\ref{sec:exactmetric}, or
from the grid with one of the local samplers of Sec.~\ref{sec:samplers})
and advances its own Runge--Kutta state.  The only synchronization within
one rank is the OpenMP barrier at the
end of the parallel region; in multi-process runs it is followed by the
MPI state exchanges described below.

\paragraph{MPI and AMR.}
The particle arrays are declared \texttt{DISTRIB=CONSTANT} (replicated on
every rank) and are checkpointable, so the particle state survives a
restart.  In a multi-process grid run the domain is partitioned among the
ranks and each rank integrates the particles owned by its local patch
(interior-box ownership with a claim margin).  A fast particle can cross
a rank face within a single Cactus step, so in this mode the step is
integrated in coordinate-time chunks: the chunk length is chosen such
that every live particle stays strictly inside its owner's interior box,
and at each chunk boundary the full state is exchanged and the owner map
is recomputed, so that a particle that has crossed into a neighbour's
box is picked up by that rank for the next chunk.  Every rank derives
the identical chunk boundaries from the identical replicated state, so
agreeing on the chunking needs no extra communication.  After the last
chunk and the loss policy, an \texttt{MPI\_Allgather} exchanges the
per-particle state, so the replicated arrays are bit-identical on every
rank.  The centred stencil is made available across rank faces
explicitly: each rank keeps a halo-padded copy of the ten
\texttt{ADMBase} fields (all three time levels) of its own patch, five
cells wide, filled by an all-pairs MPI exchange at the start of each
integration batch.  The exchange covers edges and corners as well as
faces: the three directions are exchanged in turn, and for each
direction already handled the transverse range of the later ones is
widened by the halo width, so a cell lying in the halo band of two or
three directions at once is filled as well.  Exchanging faces only
leaves those cells uninitialised, and a node Christoffel built from a
window containing one silently falls back to the analytic assembly,
which is what made the trajectory depend on the rank count.  The
chunk-boundary constraint is likewise applied even to a face the
particle sits exactly on (the distance to the face is floored at
\texttt{FACE\_EPS}); otherwise such a particle takes one unbounded
chunk and flies across the neighbouring subdomain before being dropped.
The fit-history flags travel with the particle state at each chunk
boundary, so a particle handed to a new owner refits rather than
reusing its predecessor's cached polynomials.  For the grid-fit sampler
the scheme is therefore rank-count invariant to roundoff: the same par
file run at 1, 2 and 4 ranks gives particle dumps that agree to
$1.1\times10^{-8}$ in the worst column and exactly at two ranks,
including for particles whose stencil straddles a rank face.  Exact
agreement is not expected: the chunk boundaries follow the
decomposition, so a particle is handed to its new owner at a different
point along its path and the floating-point summation order changes.
The testsuite enforces the agreement: the four-process test
\texttt{conv100\_local27} and its single-process twin
\texttt{conv100\_local27\_1p} are both compared against one and the
same reference dump, and the thorn also ships a 1-, 2- and 4-rank
agreement script with a $10^{-6}$ tolerance,
\texttt{test/check\_mpi\_identity.sh}.
The pointwise samplers instead interpolate from the rank's own patch
through \texttt{CCTK\_InterpLocalUniform}, whose stencil is clamped at
the patch boundary; there the same par file at one and at four ranks
agrees to about $3\times10^{-13}$ relative rather than bit-for-bit, and
no bit-identity is claimed for those back-ends.  It currently
supports up to four ranks (a $2\times2$ decomposition in $y$ and $z$, one
patch per rank per level).  In multi-process \texttt{Exact} runs the
analytic metric is valid everywhere, so rank~0 integrates all particles
and writes the dump.  The integration is scheduled at
\texttt{CCTK\_POSTSTEP} in Carpet's level mode: each call acquires the
calling level's grid pointers and geometry into a per-level descriptor,
and the particles are integrated once per call batch --- under Carpet's
time subcycling, on the finest-level call, with the step given by the
elapsed coordinate time --- so the integrator works unchanged in an AMR
evolution.  Each particle is integrated on the finest refinement level
whose safe region contains it (up to 16 levels,
\texttt{geodesic\_amr\_levels}, default 3), and MultiPatch
(Thornburg~\citep{thornburg2004}) block systems are supported: a
particle that crosses a block boundary has its position and spatial
velocity transformed by the block Jacobians.

\paragraph{Particle lifecycle.}  What happens to a particle that leaves
the computational domain or whose four-velocity fails the normalization
check (no future-directed root of Eq.~(\ref{eq:renorm})) is selected at
run time with \texttt{particle\_loss\_policy}: \texttt{drop} (the
default) marks the particle state as NaN and stops tracking it,
\texttt{flag} freezes it at its last valid state, and \texttt{respawn}
re-places it at a uniformly random position inside the safe box, the grid
domain inset by seven cells, subject to the rejections
$r < \mathrm{excised\_radius}+1.5$ (inside the exclusion sphere),
$r > \mathrm{initial\_radius}$ (outside the initial particle shell), and
$|z| < \mathrm{particle\_middle\_sp}$ (within
\texttt{particle\_middle\_sp} of the midplane, default $1.0$), with
$u^i = 0$ and a $u^t$ placeholder; the per-step
re-solve of
Eq.~(\ref{eq:renorm}) then normalizes the four-velocity on the next step.
Every loss is reported to the log with the particle index.  The scaling
run of Sec.~\ref{sec:scaling} uses \texttt{respawn}, which keeps the
particle count constant over long runs and prevents a single lost
particle from poisoning the statistics.

\paragraph{Diagnostics.}  With \texttt{particle\_dump\_every = k > 0} the
thorn writes the full particle state\\
$(t,x,y,z,u^t,u^x,u^y,u^z,\tau)$ to the file named by the parameter
\texttt{particle\_dump\_file} (default
\texttt{geodesic\_particles.txt}): the file begins with the initial state
($t=0$) recorded before the first integration step, followed by one
block of rows every $k$ iterations (rank~0 only); this file is the data
source for all verification figures below.
\texttt{rverbose = yes} prints the particle state at every evaluation of
the ODE right-hand side to the log for debugging.

\section{Verification}\label{sec:verification}

\subsection{Test orbits}

All verification runs use $M=2$, $a=0.5$ (the values of the \texttt{KerrSchild}
initial data \citep{moreno2002}) and a 3+1 grid with Cartesian cells of size $\Delta x$.  Three
bound orbits are used (Table~\ref{tab:orbits}), each generated and validated
independently in Python before the thorn run:

\begin{table}[h]
\centering
\caption{Test orbits used for verification ($M=2$, $a=0.5$).  $u^t$ and
$u^x,u^y,u^z$ are the four-velocity components in the par file.
$\Delta E/E$ is the value measured by the independent Python reference over
the same interval.}\label{tab:orbits}
\resizebox{\textwidth}{!}{%
\begin{tabular}{lccccc}
\toprule
orbit & initial $(x,y,z)$ & $u^t$ & $u^x,u^y,u^z$ & radial range & ref. $\Delta E/E$ \\
\midrule
P0 circular ($r=10$)   & $(10,0,0)$  & $1.53428$ & $(0,\,0.67238,\,0)$ & $r=10$ & $1.2\times10^{-9}$ \\
P1 tilted              & $(-13.448,\,1.180,\,0)$ & $1.32221$ & $(-0.01827,\,-0.50381,\,0.1)$ & $[12.86,\,16.88]$ & $5.0\times10^{-9}$ \\
P2 eccentric           & $(15,0,0)$  & $1.27607$ & $(0,\,0.45176,\,0)$ & $[12.03,\,15]$ & $5.3\times10^{-9}$ \\
\bottomrule
\end{tabular}}
\end{table}

P0 is a circular orbit of radius $r=10$ (azimuthal period
$T_\phi = 143.4\,M$); P2 is an eccentric orbit with apocenter $r_a = 15$
and pericenter $r_p = 12.03$; P1 is P2's orbit given a small vertical kick
at one point, tilting it out of the equatorial plane.
Figure~\ref{fig:orbits3d} shows the three orbits as integrated by the thorn
(Exact mode, $T=1000\,M$): P0 closes on itself in the equatorial plane,
while the rosette patterns of P1 and P2 reflect the relativistic periastron
advance of the Kerr spacetime.

\begin{figure}[h]
\centering
\maybefig{figs/fig_3orbits_3d.png}{0.85\textwidth}
\caption{The three test orbits of Table~\ref{tab:orbits} in the Kerr
spacetime ($M=2$, $a=0.5$), integrated by the thorn in Exact mode over
$T=1000\,M$ (100 snapshots).  P0, a circular orbit of radius $r=10$ in the
equatorial plane (blue); P1, P2's orbit given a small vertical kick,
precessing out of the plane (red); P2, an eccentric orbit with
$r_a = 15$, $r_p = 12.03$ (green).}
\label{fig:orbits3d}
\end{figure}

\subsection{Conservation of $E$, $L_z$ and the norm}

We integrate the three orbits simultaneously with \texttt{Exact = yes} for
$T = 1000\,M$ ($10^4$ Cactus steps of $\Delta t = 0.1\,M$, i.e.\ about seven
azimuthal periods of P0), on a $100^3$ grid.  Figure~\ref{fig:conservation}
shows the relative drift of $E$ and $L_z$ and the norm residual
$|g_{\mu\nu}u^\mu u^\nu + 1|$ for all three orbits.  The relative energy
drift stays at the level of the independent reference integration
($\le \num{maxdE}$) and the norm is held at the roundoff level
($\le \num{maxnorm}$) throughout, confirming that the RKF45 integrator with
per-step renormalization reproduces the geodesic equation accurately over
many orbital periods.

\begin{figure}[h]
\centering
\maybefig{figs/fig_conservation.pdf}{\textwidth}
\caption{Conservation over $T=1000\,M$ (Exact mode): relative drift of the
energy $E$ (left) and axial angular momentum $L_z$ (middle), and the norm
residual $|g_{\mu\nu}u^\mu u^\nu+1|$ (right), for the three orbits of
Table~\ref{tab:orbits}.  All quantities are conserved at the
$\lesssim 10^{-\num{consexp}}$ level.}
\label{fig:conservation}
\end{figure}

\subsection{Grid convergence of the interpolation mode}\label{sec:convergence}

To quantify the accuracy of the default \texttt{local27} sampler, we track
the P0 circular orbit ($M = 2$, $a = 0.5$, $r = 10$, azimuthal period
$T_\phi \simeq 143.4\,M$) for just over one full azimuthal period
($t = 150$) on grids of $50^3$, $100^3$ and $200^3$ cells
($\Delta x = 1.0$, $0.5$, $0.25$) with \texttt{Exact = no}, and compare
the final particle position with an independent Python RK4 integration of
the exact metric (step $0.02\,M$).  Figure~\ref{fig:convergence} shows
the orbit traces and the final-position error versus grid spacing: the
error decreases from $2.3\times10^{-2}$ to $4.4\times10^{-4}$ and
$9.5\times10^{-6}$ on the two refinements (factors of $52$ and $46$), a
measured local convergence order of \num{convorder}, and the exact-metric
run agrees with the reference to $\num{exacterr}$, which is the level at
which two independent integrators differ.  The quartic fit used by
\texttt{local27} is formally fifth-order accurate in space; the measured
rate of \num{convorder} on this orbit is consistent with that formal
order, and we report the measured rate as such.

\begin{figure}[H]
\centering
\maybefig{figs/fig_convergence.pdf}{\textwidth}
\caption{Grid convergence of the P0 circular orbit
($M = 2$, $a = 0.5$, $r = 10$; $t = 0\ldots 150$, just over one full
period).  Left:
orbit traces in the $x$--$y$ plane for the interpolation runs
($\Delta x = 1.0, 0.5, 0.25$) and the exact-metric run, compared with the
independent Python reference (solid).  Right: final-position error versus
grid spacing (markers), with a reference line of the measured order
\num{convorder} (dotted) and the exact-metric error level (dotted
horizontal line).}
\label{fig:convergence}
\end{figure}

Table~\ref{tab:kx100} gives the maximum position error of each grid
sampler against the \texttt{Exact = yes} trajectory of the same orbit.

\begin{table}[h]
\centering
\caption{Maximum position error against the \texttt{Exact = yes}
trajectory: P0 circular orbit ($M = 2$, $a = 0.5$, $r = 10$) on the
static Kerr--Schild field of a $100^3$ grid ($\Delta x = 0.5\,M$),
$\Delta t = 0.1\,M$, single process, particle started at $t = 0$.}\label{tab:kx100}
\begin{tabular}{lccc}
\toprule
metric evaluation & $t = 2\,M$ & $t = 10\,M$ & $t = 20\,M$ \\
\midrule
\texttt{l27}     & $3.74\times10^{-6}$ & $4.42\times10^{-5}$ & $9.97\times10^{-5}$ \\
\texttt{aei}     & $3.68\times10^{-7}$ & $6.29\times10^{-6}$ & $2.97\times10^{-5}$ \\
\texttt{local27} & $1.60\times10^{-7}$ & $2.56\times10^{-7}$ & $1.97\times10^{-7}$ \\
\bottomrule
\end{tabular}
\end{table}

\subsection{Comparison of the grid samplers on a precessing orbit}\label{sec:sampler-compare}

To compare the grid samplers on an orbit closer to production use,
we track $100$ particles that all move on the same bound equatorial Kerr
ellipse ($M=2$, $a=0.5$, pericentre $r_p = 16$, apocentre $r_a = 64$,
$E = 0.9761517$, $L_z = 8.1632196$), spread over the orbit, on the static
Kerr--Schild field written onto a $144^3$ base grid ($\Delta x = 1$) with
three nested refinement boxes ($\Delta x = 0.5$, $0.25$, $0.125$ out to
radii $36$, $18$, $9$).  The ellipse precesses by $117.5^\circ$
($0.6528\pi$) per radial period ($T_r = 1349.7$), so the orbit stays
inside the computational domain ($16 \le r \le 64$; domain half-size
$72$) for the whole run, $t = 0 \ldots 3200\,M$ (about $2.37$ radial
periods); the particle state is dumped every $10$ time units, and a
particle is counted as lost
when it leaves the safe region of its grid patch or fails an ODE step
(the thorn's drop policy).  The six runs are \texttt{Exact = yes}
(reference, single process, 20 OpenMP threads),
\texttt{metric\_interp = aei} and \texttt{metric\_interp = carpet}
(each single process, 20 threads; both evaluate the same AEI order-4
pointwise Lagrange interpolant and differ only in the
data-acquisition path), and \texttt{metric\_interp = local27} in three
configurations: the default four-process mode (four processes, 20
threads each), the explicit \texttt{local27} four-process mode (four
processes, 20 threads each), and a single-process control (one process,
20 threads).

\begin{table}[h]
\centering
\caption{Comparison of the grid samplers on the precessing
ellipse of Sec.~\ref{sec:sampler-compare} ($100$ particles,
$t = 0 \ldots 3200\,M$; no particle lost in any run).
$\max|\Delta x|$: largest per-particle coordinate deviation from the
\texttt{exact} reference, over the run.
$\Delta E$, $\Delta L_z$: maximum and mean absolute deviation of the
conserved quantities from their initial values, over the run.
"p0": radial position of particle 0 at the end of the run.}\label{tab:samplers}
\resizebox{\textwidth}{!}{%
\begin{tabular}{llccccccl}
\toprule
metric evaluation & run & lost & survivors & $\max|\Delta x|$ vs ref & $\max|\Delta E|$ / mean & $\max|\Delta L_z|$ / mean & p0 \\
\midrule
exact (reference)  & 1 process, 20 threads     & 0 & 100 & --- & --- & --- & $r=32.47$ (alive) \\
aei                & 1 process, 20 threads     & 0 & 100 & $2.09\times10^{-2}$ & $1.14\times10^{-6}$ / $1.38\times10^{-7}$  & $6.53\times10^{-5}$ / $1.97\times10^{-5}$  & $r=32.47$ (alive) \\
carpet             & 1 process                 & 0 & 100 & (identical to \texttt{aei}) & (identical to \texttt{aei}) & (identical to \texttt{aei}) & $r=32.47$ (alive) \\
local27 (explicit) & 4 processes, 20 threads ea. & 0 & 100 & $2.43\times10^{-1}$ & $1.32\times10^{-5}$ / $2.26\times10^{-6}$  & $6.16\times10^{-4}$ / $2.87\times10^{-4}$  & $r=32.38$ (alive) \\
local27 (default)  & 4 processes, 20 threads ea. & 0 & 100 & (identical to explicit) & (identical to explicit) & (identical to explicit) & $r=32.38$ (alive) \\
local27 (control)  & 1 process, 20 threads     & 0 & 100 & $2.28\times10^{-1}$ & $1.07\times10^{-5}$ / $2.01\times10^{-6}$  & $6.92\times10^{-4}$ / $3.44\times10^{-4}$  & $r=32.44$ (alive) \\
\bottomrule
\end{tabular}}
\end{table}

\begin{figure}[H]
\centering
\maybefig{figs/fig_interp_compare.png}{0.9\textwidth}
\caption{Particle 0 of the precessing-ellipse comparison
(Sec.~\ref{sec:sampler-compare}) in the $x$--$y$ plane: exact reference
(black), the two \texttt{local27} four-process runs (blue solid and gray
dashed, bit-identical to each other), the \texttt{local27}
single-process control (teal dashed), \texttt{aei} (red) and
\texttt{carpet} (green; bit-identical to \texttt{aei} in this test).
Circles: apocentre ($r=64$), pericentre ($r=16$) and the excision
sphere ($r=4$, dashed).}\label{fig:interp-compare}
\end{figure}

\begin{figure}[H]
\centering
\maybefig{figs/fig_interp_compare_p0.png}{\textwidth}
\caption{Radial coordinate $r(t)$ of particle 0 for the runs of
Table~\ref{tab:samplers}.  Dashed: pericentre ($r=16$) and apocentre
($r=64$) of the reference ellipse.}
\label{fig:interp-compare-p0}
\end{figure}

No particle is lost in any of the six runs: all $100$ particles
complete the full $t = 0 \ldots 3200$ run without leaving the safe
region of their grid patch or failing an ODE step.  The deviation from
the exact-metric reference separates the samplers cleanly
(Table~\ref{tab:samplers},
Figure~\ref{fig:interp-compare-p0}): the two AEI order-4 runs
(\texttt{aei}, \texttt{carpet}; bit-identical to each other) track the
reference to $2.1\times10^{-2}$, and both \texttt{local27}
configurations stay within $2.3$--$2.4\times10^{-1}$.  In every run the
dominant error is the grid's own coherent discretization of the exact
Kerr--Schild background, which the three samplers reproduce at the
$10^{-2}\ldots10^{-1}$ level; the node-exact \texttt{local27} quartic
fit carries that background error further than the AEI order-4
interpolant ($2.3\times10^{-1}$ against $2.1\times10^{-2}$).  The
four-process and single-process \texttt{local27} runs now differ by no
more than this background error: on this mesh one and four processes
agree to $2.0\times10^{-1}$ in the coordinates at $t = 3200\,M$, and
their $t = 0$ frames are bit-identical.

Before the fixes described in Sec.~\ref{sec:impl} the same
one-versus-four-process comparison was $O(1)$ --- $2.35$ in $x$ --- and
the cause was not lost precision but a different metric: the owner map
used to pick a refinement level together with an owning rank by distance
to that rank's published box, a level-0 interior box ends exactly at the
rank face (zero gap) while a refined level's claim box is inset by seven
cells, so particles on the $z = 0$ plane --- where this $2\times2$
decomposition puts its faces --- were sampled at $\Delta x = 1$ with four
processes and at $\Delta x = 1/2$ with one.  Selecting the level first,
from the union of the patch extents of all ranks on that level (the same
criterion \texttt{pick\_level} applies in a single process), removes the
discrepancy at $t = 0$ entirely.

What remains is a chunking artefact rather than a physical one.  Under
AMR with \texttt{Geodesic::step} $\gg \Delta t$ the chunk boundaries fall
on rank faces; a single process has no internal faces, so the RKF45
segmentation and the points at which the owner level is re-evaluated
differ, and the residual grows roughly linearly with coordinate time
($7.9\times10^{-6}$ at $t = 10$, $6.0\times10^{-2}$ at $t = 1600$,
$1.6\times10^{-1}$ at $t = 3200$).  At the shorter AMR integration of the
testsuite (\texttt{extra/ellipse\_orbit\_amr.par}, four levels,
$t = 25.6$) the one-, two- and four-process dumps are still
bit-identical.  On a unigrid mesh, where no internal
faces exist, the same comparison is at roundoff level:
\texttt{conv100\_local27} agrees to $1.1\times10^{-8}$ across 1, 2 and 4
ranks, and there the deviation is bounded and oscillates with the orbit
instead of growing (Sec.~\ref{sec:impl}).  This has no adverse
consequence in production use: in a full \texttt{GRHydro} evolution, or
any grid-carried spacetime, the grid field \emph{is} the true metric,
and the single-process runs show that the tracker follows it faithfully.

\paragraph{Run notes.}
The default-mode four-process run and the explicit \texttt{local27}
four-process run are bit-identical to each other (identical mode,
different par spelling), confirming the default path of
Sec.~\ref{sec:samplers}.  The \texttt{aei} and \texttt{carpet} rows are
taken from single-process runs.  Both back-ends also run at four
processes on this mesh, and do so without loss: their $t = 0$ frames are
bit-identical to the single-process ones and all $100$ particles survive
to $t = 3200\,M$, but the chunking artefact described above is not
specific to the fitted sampler --- four- and single-process
\texttt{aei} separate by $9.5\times10^{-2}$ over the run, against
$2.1\times10^{-2}$ (single process) and $8.5\times10^{-2}$ (four
processes) from the exact reference ---
so the table quotes the single-process \texttt{aei} and \texttt{carpet}
runs, whose deviation from the reference is the smaller and more cleanly
defined one.  It is the \texttt{l27} path that must be run
single-process (Sec.~\ref{sec:samplers}).

\subsection{Geodesic orbits on an AMR mesh}\label{sec:amr-orbits}

The grid-mode samplers read the metric from the refinement level that
contains the particle, which makes them usable on Carpet's AMR meshes
without modification.  We verify this on the same circular orbit as the
unigrid reference \texttt{conv100.par} (Kerr--Schild, $M=2$, $a=0.5$,
particle on the $r=10$ equatorial circle, integrated to $t=150.1$,
$\sim 1.05$ orbital periods, $T\simeq 143$), now evolved on an AMR mesh
with three levels of resolution: a $50^3$ base grid ($\Delta x=1$)
refined twice (factor 2) by a box centred on the origin (half-width $16$
cells on the first refinement level, then $8$), so that the finest level
has $\Delta x=1/4$.  The orbiting particle lies on the first refinement
level ($\Delta x=1/2$) for its whole path --- the same local resolution
as the $100^3$ unigrid reference --- and crosses onto the second level
($\Delta x=1/4$) near the box diagonals.  The particle integrator is
otherwise unchanged (RKF45, same seed and initial data as the reference
run), so the only differences between the two set-ups are the grid
structure and the driver's 4:1 subcycling.

\begin{table}[h]
\centering
\caption{AMR (3 levels) vs.\ unigrid reference (conv100): maxima over
the 16 dump times $t=0.1,10.1,\ldots,150.1$ of the Cartesian position
difference $\max_i|\Delta x^i|$, the radial difference $|\Delta r|$, and
the orbital phase difference $|\Delta\phi|$.  The reference run leaves
\texttt{metric\_interp} unset and so uses the default \texttt{local27};
each row is that sampler's AMR-to-unigrid difference, not a ranking of
sampler accuracy.}\label{tab:amr-orbits}
\begin{tabular}{lccc}
\toprule
sampler & $\max_i|\Delta x^i|$ & $\max|\Delta r|$ & $\max|\Delta\phi|$ (rad) \\
\midrule
local27           & $1.2\times10^{-1}$ & $4.3\times10^{-2}$ & $1.3\times10^{-2}$ \\
aei               & $1.4\times10^{-1}$ & $5.0\times10^{-2}$ & $1.6\times10^{-2}$ \\
carpet            & $3.4\times10^{-1}$ & $2.3\times10^{-1}$ & $4.2\times10^{-2}$ \\
\bottomrule
\end{tabular}
\end{table}

All three samplers keep the particle on a stable orbit throughout the
integration (no blow-up, no lost particle).  \texttt{local27} and \texttt{aei}
stay within $5\times10^{-2}$ ($0.5\%$) of the reference in radius and
$1.7\times10^{-2}$ rad ($\sim1^\circ$) in phase; \texttt{carpet} drifts most
on this mesh, $2.3\times10^{-1}$ ($2.3\%$) in radius.  The reference run
(\texttt{conv100.par}) leaves \texttt{metric\_interp} unset and so uses the
default \texttt{local27}: the \texttt{local27} row measures the effect of the
mesh alone, while the \texttt{aei} and \texttt{carpet} rows include a sampler
change as well, so this table does not rank the samplers by accuracy (the accuracy
comparison is Sec.~\ref{sec:sampler-compare}).  At the first dump after the mesh
settles ($t=10.1$) all three deviations are already at their small-$t$ level
($\lesssim5\times10^{-3}$ in position): no transient is associated with the
refinement, the pointwise samplers resolve the metric correctly across level
boundaries, and the geodesic equation integrates in a Carpet AMR evolution (with
its 4:1 subcycling) without modifying the integrator.

\subsection{Robustness on a two-puncture merger background}\label{sec:gw150914}

The checks above all use a single-black-hole Kerr--Schild field, where the
dominant error source is the grid's own discretization.  A closer production use
case is the field of a binary black hole, whose geodesic structure is itself of
observational interest \citep{chernov2024}.  We therefore place the tracker on
the two-puncture initial data of the \texttt{GW150914} binary (the first
directly detected gravitational-wave signal \citep{abbott2016gw150914};
component masses $36$ and $29\,M_\odot$, stored in the par files as
$m_+ = 36/65$ and $m_- = 29/65$, i.e.\ normalized to total mass $M = 1$), built
by \texttt{TwoPunctures} with the GW150914 spin and kick parameters
(\texttt{par\_b} = 5, punctures at $x = +4.46$ and $-5.54$), on a Cartesian
domain of $\pm10.5$, evolved with \texttt{ML\_BSSN} (8th-order SBP differences
$+$ MoL RK4, 1+log slicing, $\Gamma$-driver shift).  There is no analytic
reference orbit on this background, so this section is about stability
thresholds and the cost of resolution and timestep, not about absolute
accuracy.

An important caveat: the raw \texttt{TwoPunctures} initial data is not the
merger metric itself.  Once the field is evolved with \texttt{ML\_BSSN}, it
relaxes and differs significantly from the frozen initial geometry.  We
therefore make the time-dependent, \texttt{ML\_BSSN}-evolved background the main
robustness test, and treat the frozen initial data only as a secondary check.

Common protocol: a $100$-particle equatorial ring at $r = 8$ (azimuthal spacing
$3.6^\circ$, strong inward velocity; particle 0 sits at $(8,0,0)$ with
$u^t = 1.774$, $u^x = -0.960$, $u^y = 0.863$), released at $t = 1.60$ (after the
initial junk transient, which otherwise changes the drop statistics), dumped
every $0.05$ in coordinate time, integrated to $t = 14.05$; $4$ MPI processes
$\times\,21$ threads.  A binary has two punctures, so an excision sphere
centered on the origin is meaningless and can interfere with the computation: the
current protocol sets \texttt{excised\_radius}$ = 0$.  The older runs in the
table below still carry $2.0$, but that sphere never fired --- dropped particles
get no closer than $4.24$ to the origin (they fall into the punctures at
$x = +4.46$ and $-5.54$) and the smallest radius reached by any particle is
$2.33$--$2.38$ --- so the survival counts remain comparable.

\paragraph{Resolution and timestep.}
Table~\ref{tab:tpgrid} collects six evolved-background runs under this protocol
plus four coarse-grid probes.  The ``background outcome'' column records whether
the background field itself integrates through $t = 14.05$.

\begin{table}[h]
\centering
\caption{$r = 8$ ring, $100$ particles, released at $t = 1.60$, dumped every
$0.05$ in coordinate time, integrated to $t = 14.05$ ($2248$ steps at
$\Delta t = 0.00625$ on the unigrid runs, $1124$ finest-level steps at
$\Delta t = 0.0125$ on the AMR run).  \emph{wall}: measured wall-clock time of
one run ($4\times21$ threads); \emph{lag}: median of (grid time
$-$ the particle's own integrated coordinate time) at $t = 14.05$, or at the last
frame with $\ge50\%$ of the particles alive for the coarse probes;
\emph{alive}: particles not inside a puncture at the end of the run.}
\label{tab:tpgrid}
\resizebox{\textwidth}{!}{%
\begin{tabular}{llcccc}
\toprule
grid & $\Delta x$ & $\Delta t$ & wall (s) & lag & alive / background outcome \\
\midrule
AMR $1/8$ base $+\,1/16$ box ($r=8$) & $0.125/0.0625$ & $0.025/0.0125$ & $51{,}818$ & $+0.013$ & $72$ / stable \\
\quad same grid, \texttt{aei} sampler & as above & as above & $54{,}782$ & $+0.033$ & $76$ / stable \\
unigrid $1.8/16$ ($187^3$) & $0.1123$ & $0.00625$ & $28{,}545$ & $0$ & $74$ / stable \\
\quad same grid, \texttt{aei} sampler & as above & as above & $25{,}014$ & $0$ & $81$ / stable \\
unigrid $1.5/16$ ($224^3$) & $0.09375$ & $0.00625$ & $47{,}340$ & $0$ & $74$ / stable \\
\quad same grid, \texttt{aei} sampler & as above & as above & $44{,}673$ & $0$ & $78$ / stable \\
\midrule
unigrid $1/4$ ($74^3$) & $0.25$ & $0.05$ & --- & $+0.155$ & ring lost before $11.60$ \\
unigrid $1/8$ ($148^3$) & $0.125$ & $0.025$ & --- & $+0.049$ & ring lost before $11.60$ \\
unigrid $1/8$ ($168^3$) & $0.125$ & $0.0125$ & --- & $0$ & background blows up at $t=10.85$ (shift $\sim10^{246}$) \\
unigrid $1/8$ ($168^3$) & $0.125$ & $0.00625$ & --- & $0$ & background blows up at $t=10.65$ (shift $\sim10^{18}$) \\
\bottomrule
\end{tabular}}
\end{table}

Three conclusions.  \emph{First, what limits the run is the resolution of the
background, not the particles.}  On the $\Delta x = 0.125$ grid the shift at
$t = 10.8625$ is still regular, $(-0.091, -0.032, -5\times10^{-4})$, and one
timestep later ($\Delta t = 0.0125$) it is $\sim10^{246}$: no precursor, no
\texttt{NaN}, a grid-scale mode at the punctures being amplified.  Halving
$\Delta t$ moves the blow-up by only $\sim0.2\,M$.  The threshold lies between
$\Delta x = 0.125$ (fails) and $0.1123$ (succeeds), and has nothing to do with
the geodesic integrator.  \emph{Second, the particle-clock lag is a
$\Delta t$ effect.}  At fixed $\Delta x$ it falls roughly in proportion to
$\Delta t$ ($+0.155$ at $0.05$, $+0.049$ at $0.025$, $0$ at $0.0125$ and below),
and it is $+0.013$/$+0.033$ on the AMR runs.  The mechanism is the
\texttt{max\_iter\_cap} bound on the proper-time substeps inside one Cactus
step: when it binds, the particle never reaches the target coordinate time, and a
coarser grid (larger interpolation error) pushes RKF45 to smaller internal steps,
making the bound more likely to bind.
\emph{Third, under a fixed $\Delta x\,\Delta t$ constraint, shrinking
$\Delta t$ is far cheaper than shrinking $\Delta x$:} the step count grows
linearly while the cell count grows by a factor of $8$ at the same volume.  As
long as the background produces no \texttt{NaN}, refining the timestep is the
preferred knob.

\paragraph{Particle clock and sampler consistency.}
The dump now carries the grid coordinate time (\texttt{cctk\_time}) as a column,
so the lag above can be checked directly.  When two runs are compared, both
trajectories must first be interpolated to a common coordinate time: pairing by
dump index turns the lag into an along-track offset: on the AMR \texttt{aei} row
the median position difference at $t = 14.05$ inflates from $0.0043$ to $0.0159$
($3.7\times$) and the rms median from $0.0039$ to $0.0063$; on unigrid runs the
lag is zero and the two pairings agree bit for bit.
Table~\ref{tab:tpdiff} gives the interpolated comparison, with the unigrid
$1.5/16$ \texttt{local27} run as reference.

\begin{table}[h]
\centering
\caption{Relative to the unigrid $1.5/16$ \texttt{local27} run:
$d_{\mathrm{pos}}$ is the median position difference at the stated coordinate
times over particles inside the trusted region in both runs;
\emph{rms\_med}/\emph{rms\_max} are the median and maximum over particles of
the rms position difference accumulated over the trusted part of the trajectory;
$d\tau$ is the median $|\Delta\tau|$ at each particle's last trusted frame;
$\Delta t_{\mathrm{drop}}$ is the median difference in drop time over particles
captured in both runs, and $n_{\mathrm{cap}}$ is that number; alive is the count
at the end of the run.  The trusted region is $r > 3$ from the origin and more
than $1$ from either puncture, so nothing is measured after a particle has fallen
in.  The sample is all $100$ particles minus two chaotic outliers (\#53, \#54),
i.e.\ $98$; removing them changes only rms\_max (by $2.9$--$7.5\times$), the
median quantities being unchanged in the fourth decimal place.}
\label{tab:tpdiff}
\resizebox{\textwidth}{!}{%
\begin{tabular}{lccccccc}
\toprule
run & $d_{\mathrm{pos}}$ @$t=6/10/14$ & rms\_med & rms\_max & $d\tau$ & $\Delta t_{\mathrm{drop}}$ & $n_{\mathrm{cap}}$ & alive \\
\midrule
AMR $1/8+1/16$, \texttt{local27} & $0.0041/0.0041/0.0050$ & $0.0040$ & $0.0109$ & $0.0052$ & $0.30$ & $26$ & $72$ \\
AMR $1/8+1/16$, \texttt{aei}     & $0.0040/0.0037/0.0043$ & $0.0039$ & $0.0099$ & $0.0046$ & $0.20$ & $24$ & $76$ \\
unigrid $1.8/16$, \texttt{local27} & $0.0000/0.0005/0.0018$ & $0.0009$ & $0.0207$ & $0.0008$ & $0.25$ & $25$ & $74$ \\
unigrid $1.8/16$, \texttt{aei}      & $0.0001/0.0007/0.0018$ & $0.0009$ & $0.0207$ & $0.0008$ & $1.25$ & $19$ & $81$ \\
unigrid $1.5/16$, \texttt{aei}      & $0.0000/0.0002/0.0003$ & $0.0001$ & $0.0010$ & $0.0002$ & $0.45$ & $22$ & $78$ \\
unigrid $1.5/16$, \texttt{local27} (reference) & $0$ & $0$ & $0$ & $0$ & $0$ & $0$ & $74$ \\
\bottomrule
\end{tabular}}
\end{table}

Switching sampler on the same $1.5/16$ grid
(\texttt{local27}$\to$\texttt{aei}) changes positions at $t = 14.05$ by a median
of $3.2\times10^{-4}$, a 90th percentile of $8.6\times10^{-4}$ and a worst case
of $8.2\times10^{-3}$; switching resolution from $1.8/16$ to $1.5/16$ gives
$1.8\times10^{-3}$, $1.8\times10^{-2}$ and $5.1\times10^{-1}$.  The choice of
sampler therefore matters one to three orders of magnitude less than the
resolution.

\paragraph{Survival is not a quality metric.}
The survival count depends on the background and on the sampler, so it cannot
rank accuracy: in the $1.8/16$ \texttt{local27} run the $26$ dropped particles
were $0.15$--$0.36$ from a puncture (median $0.28$) at their last finite
snapshot, other runs give $0.08$--$0.50$, and particles that survive get as close
as $0.24$--$0.26$; on the same $1.5/16$ grid, changing sampler alone moves the
median drop time by $0.45$, the same size as the difference between resolution
levels ($0.25$).  Figures~\ref{fig:ring-u1516-local27} and
\ref{fig:ring-u1516-aei} show the trajectories on the finest grid for both
samplers, and their morphology is the same.

\begin{figure}[H]
\centering
\maybefig{figs/fig_ring_u1516_local27.png}{0.92\textwidth}
\caption{$r = 8$ ring on the unigrid $1.5/16$ background ($224^3$,
$\Delta x = 0.09375$, $\Delta t = 0.00625$), \texttt{local27} sampler,
$t = 1.60 \ldots 14.05$ ($250$ snapshots, colour = particle index); $\times$
marks the punctures located by \texttt{PunctureTracker}.  $74$ of the $100$
particles are alive at $t = 14.05$.}
\label{fig:ring-u1516-local27}
\end{figure}

\begin{figure}[H]
\centering
\maybefig{figs/fig_ring_u1516_aei.png}{0.92\textwidth}
\caption{Same background and same initial ring, with the pointwise
\texttt{aei} sampler: $78$ particles alive at $t = 14.05$, trajectory morphology
identical to Fig.~\ref{fig:ring-u1516-local27}, median per-particle position
difference $3.2\times10^{-4}$.}
\label{fig:ring-u1516-aei}
\end{figure}

\paragraph{What AMR costs.}
Each extra refinement level costs a factor of $16$ in volume, so a ``base grid
$+\,1/N$ box'' covers $(1 + 16f)$ times the grid points of the base grid, where
$f$ is the box volume as a fraction of the domain; this is a point-count
estimate, not a measured throughput.  A box covering the $r = 8$ ring has
$f = 0.44$, i.e.\ $8.1\times$, so the estimate predicts only a $2.1\times$
saving over refining the whole domain, and the measured wall times show no
saving at that $f$: the AMR run takes $51{,}818$ s against $47{,}340$ s for the
unigrid $1.5/16$ run ($+9\%$), and $54{,}782$ s against $44{,}673$ s for the
\texttt{aei} pair ($+23\%$).  On the same estimate AMR pays off only for
$f \lesssim 0.055$ (a box of half-width about $4$, a $1.9\times$ speedup).
More importantly, the binding constraint is
a property of the background: a refinement box around the particles does not
make a coarse background stable.  When the particles are not confined to a small
region of the domain, a unigrid with a suitable $\Delta t$ is both simpler and
cheaper; when they are, AMR still has an advantage.  As a functional check of
the AMR path we also ran a multi-patch mesh with seven static-box levels for
$2048$ finest-level steps ($t \simeq 0 \ldots 4.4$): all $100$ particles remain
finite throughout, with zero drops, on both \texttt{local27} and \texttt{aei},
which agree to $1.0\times10^{-1}$ in position (the discretization level of that
grid), showing that metric sampling across refinement boundaries is correct.

\paragraph{Frozen initial-data check.}
On the frozen raw initial data (a pure geometry-routing/stability check, not
representative of the merger), the far ring ($r = 10$, prograde tangential
velocity $0.32$, inward radial $0.10$) completes the full $80$-unit run on all
three back-ends (\texttt{local27}, \texttt{aei}, \texttt{carpet}; the latter two
sample through the \texttt{AEILocalInterp}/\texttt{LocalInterp} interface and
are bit-identical to each other) with zero \texttt{NaN} states, and the \texttt{local27}
trajectory agrees with the \texttt{aei}/\texttt{carpet} ones to within
$\sim6\times10^{-4}$ in position; on the $r = 8$ near ring, deterministic
\texttt{NaN} drops (\texttt{loss\_policy = drop}) set in from $t = 6$ and reach
exactly $50$ of the $100$ particles by $t = 20$, with an almost identical drop
pattern on all three back-ends: \texttt{aei} and \texttt{carpet} agree bit for
bit, and \texttt{local27} differs in only two of the fifty dropped particles.
Every drop is a physical excursion into the strong-field/excised region rather
than a numerical failure of the integrator: on a Kerr background these orbits
are the plunging geodesics whose end states are known analytically
\citep{dyson2023}.  All six frozen runs (near and far ring, three back-ends
each) are part of the thorn's automatic validation suite
(\texttt{test/test.ccl}); the \texttt{NaNChecker}-aware test comparator turns
the deterministic drop pattern of the near runs into a regression test.

\paragraph{Resolution scan on the frozen data.}
Because the two-puncture background has no analytic reference, we repeat the
near-ring run at three unigrid resolutions with the cell count held fixed at
$144^3$ and the domain shrunk in step ($\pm72$, $\pm36$, $\pm18$, i.e.\
$\Delta x = 1$, $0.5$, $0.25$), so all three runs resolve the puncture region
with the same number of cells; the integration parameters are identical in the
three runs, only the grid resolution changes.  On this frozen data
\texttt{aei} and \texttt{carpet} are bit-identical at every resolution
(Sec.~\ref{sec:sampler-compare} already showed the two to be the same
computation), so they are reported as one column.  Table~\ref{tab:tpconv}
reports the median position and velocity difference at identical snapshot times
between consecutive resolutions, over the particles alive in both runs of each
pair.  We report medians rather than maxima: chaotic trajectory divergence near
the punctures dominates the maximum and makes its apparent order meaningless.
$58$--$59$ of the $100$ particles fall into the excision region at different
snapshot times in the three resolutions; this is trajectory-level sensitivity to
the grid scale near the punctures, not a scheme-level failure.  For the order-4
pointwise Lagrange samplers (\texttt{aei}/\texttt{carpet}) the observed order
sits at the design order $4$; for \texttt{local27} it rises with resolution
($3.3\to7.0$) around the formal order $5$ of its quartic tensor-product fit
(Sec.~\ref{sec:convergence}), so that column is a trend rather than a single
order.

\begin{table}[h]
\centering
\caption{Resolution scan on the frozen \texttt{GW150914} two-puncture initial
data: $r = 8$ ring, $100$ particles, $T = 20\,M$, $\Delta x = 1 / 0.5 / 0.25$ at
fixed $144^3$ cells.  $p_{\mathrm{med}}$: observed order of the median position
difference between the $\Delta x = 0.5$ and $0.25$ runs relative to the
$\Delta x = 1$ and $0.5$ runs, over particles alive in both runs of each pair;
$p_v$: same for the velocity difference at $t = 20$.  Bucketed rows: median
position difference for particles whose minimum radius stays above $4$
(\emph{far}) or dips below (\emph{near}), with the observed order $p$ in
parentheses.  \texttt{aei} and \texttt{carpet} are bit-identical at all three
resolutions and reported as one column.}
\label{tab:tpconv}
\resizebox{\textwidth}{!}{%
\begin{tabular}{llcccc}
\toprule
sampler & quantity & $t=2$ & $t=10$ & $t=20$ & $p_v$ ($t=20$) \\
\midrule
\texttt{aei}/\texttt{carpet} & $p_{\mathrm{med}}$ & 4.22 & 4.14 & 4.33 & 5.61 \\
\texttt{local27}             & $p_{\mathrm{med}}$ & 3.26 & 4.75 & 6.95 & 5.75 \\
\midrule
\texttt{aei}/\texttt{carpet} & far ($n=85$)  & --- & $8.2\times10^{-5}$ ($4.08$) & $1.1\times10^{-3}$ ($4.66$) & --- \\
\texttt{aei}/\texttt{carpet} & near ($n=15$) & --- & $2.6\times10^{-3}$ ($3.95$) & $7.6\times10^{-3}$ ($4.08$) & --- \\
\texttt{local27}             & far ($n=86$)  & --- & $7.6\times10^{-5}$ ($5.36$) & $4.5\times10^{-4}$ ($6.07$) & --- \\
\texttt{local27}             & near ($n=14$) & --- & $9.0\times10^{-2}$ ($1.83$) & $3.4\times10^{-3}$ ($8.09$) & --- \\
\bottomrule
\end{tabular}}
\end{table}



\section{Strong scaling}\label{sec:scaling}

\begin{figure}[H]
\centering
\maybefig{figs/fig_scaling.pdf}{\textwidth}
\caption{Strong scaling with $10^4$ particles, $50^3$ grid, 512 steps,
\texttt{l27} sampler.  Left: wall time versus number of OpenMP
threads.  Right: speedup versus thread count (solid line, ideal) with the
parallel efficiency annotated at each point.}
\label{fig:scaling}
\end{figure}

The per-particle independence of the integration loop
(Sec.~\ref{sec:impl}) should translate directly into parallel speedup.  We
run a configuration with $10^4$ particles (the three orbits of
Table~\ref{tab:orbits} plus $9997$ particles with random positions inside
the safe box and zero velocity, fixed random seed), $50^3$ grid,
\texttt{Exact = no}, $512$ steps, on a two-socket node with 44 physical
cores (88 logical threads via hyperthreading), varying only
\texttt{OMP\_NUM\_THREADS} from 1 to 88.  Figure~\ref{fig:scaling} shows the
wall time and the parallel efficiency.  The speedup grows with the thread
count up to 32 threads ($8.3\times$) and then plateaus at $9.2\times$
(efficiency \num{scaleeff}\%) at 88 threads.  A per-phase timing breakdown
(Cactus \texttt{TimerReport}) attributes the plateau to two effects.
First, the surrounding Cactus framework (metric bookkeeping, boundary
conditions, NaN checks) consumes $\sim\!40$\,ms per step, nearly independent
of the thread count ($19.0$\,s total at one thread versus $22.9$\,s at 88):
about $3\%$ of the $\sim\!1.45$\,s single-thread step but about $29\%$ of
the $\sim\!0.15$\,s 88-thread step.
Second, the parallel particle region itself scales sublinearly: it shrinks
from $\sim\!1.41$\,s to $\sim\!0.104$\,s per step, a speedup of only
$13.7\times$ on 88 threads, consistent with the 44 physical (hyperthreaded)
cores, cross-NUMA stencil reads, and OpenMP scheduling overhead.

As a complementary measurement in a production setting, we ran the tracker
inside full \texttt{GRHydro} evolutions \citep{font2008} with $5000$ and with $10^4$
particles; each run shared the 88-thread node with three other concurrent
instances, i.e.\ only $\sim\!22$ threads were available per instance.  The
particle integration (\texttt{geodesics\_integrate} in the
\texttt{TimerReport}) consumed $49.7$\,s of $2860.5$\,s total wall time
($1.7\%$) in the $5000$-particle run and $84.1$\,s of $2749.1$\,s
($\num{hydrofrac}\%$) in the $10^4$-particle run: the dominant costs remain
the hydrodynamic reconstruction ($\sim\!27\%$) and the Riemann solves
($\sim\!10\%$), and doubling the particle count grows the integration cost
by a factor of only $1.7$.  Because these fractions were measured under
contention from the co-running instances, the overhead in a dedicated run
is smaller still: tracking test particles alongside an evolution is
essentially free compared to the evolution itself.

\section{Conclusion and availability}\label{sec:concl}

We have presented \texttt{Geodesic}, an Einstein Toolkit thorn that
integrates test-particle geodesics in Kerr--Schild Cartesian coordinates.
The thorn supports an exact analytic Kerr--Schild metric (grid-independent,
verified to be a vacuum solution to beyond machine precision) and a grid
mode
with four local samplers selected by one parameter (\texttt{l27},
\texttt{local27}, \texttt{aei}, \texttt{carpet}), each with its time
handling built in.  The grid mode
assumes no particular analytic form of the metric; for stationary grids
$\partial_t g_{\mu\nu} = 0$ is detected automatically, and for
time-dependent grids the time derivative is recovered from the stored
\texttt{ADMBase} time levels (Sec.~\ref{sec:metric}): with
\texttt{Exact = no}, the thorn tracks geodesics in whatever spacetime
the grid carries, stationary or evolving (Sec.~\ref{sec:metric}) ---
for example the two-puncture \texttt{GW150914} binary background of
Sec.~\ref{sec:gw150914}.  Multi-process runs distribute the particles
over MPI ranks with an owner-computes scheme that keeps the replicated
state bit-identical on every rank; AMR meshes with up to 16 refinement
levels and MultiPatch block systems are supported (Sec.~\ref{sec:impl}).  In that mode each particle is an
independent unit of work, and an OpenMP parallelization of the particle loop
scales to $10^4$ particles, reaching a speedup of
$9.2\times$ (\num{scaleeff}\% efficiency) on the 88 logical threads of a
two-socket, 44-core hyperthreaded node --- i.e.\ the loop is
communication-free but bandwidth-limited, not load-limited
(Sec.~\ref{sec:scaling}).  Verification against an
independent reference implementation shows conservation of $E$, $L_z$ and the
norm to $\sim10^{-\num{consexp}}$ over $10^3\,M$, and a measured local
convergence order of \num{convorder} for the default sampler.  On a precessing
orbit of $100$ particles the AEI order-4 path agrees with the exact-metric
reference at $2.1\times10^{-2}$ and \texttt{local27} at
$2.3$--$2.4\times10^{-1}$, with one and four processes agreeing to
$2.0\times10^{-1}$ on that mesh
(Sec.~\ref{sec:sampler-compare}); on a 3-level AMR mesh the grid samplers
reproduce the unigrid reference to $4.3\times10^{-2}$--$2.3\times10^{-1}$ in
radius (Sec.~\ref{sec:amr-orbits}).  On the \texttt{GW150914} two-puncture
merger background geodesics integrate through $t = 14.05\,M$ whenever the
background itself is stable, and the cost is set mainly by the resolution of the
background, $\Delta t$ being the cheaper knob (Sec.~\ref{sec:gw150914}).

The thorn is released under the GPLv2+ license together with the
configuration files used in this work and with the source of this manuscript
and its figures (\texttt{paper/} and \texttt{paper\_zh/}), at
\url{https://github.com/wylyhzh/Geodesic}.  Of the 46 parameter files in \texttt{doc/par/} of that
repository, 32 are the runs quoted above: the conservation and
grid-convergence runs of Secs.~\ref{sec:convergence}, the sampler
comparison of Sec.~\ref{sec:sampler-compare}, the AMR orbits of
Sec.~\ref{sec:amr-orbits}, the \texttt{GW150914} runs of
Sec.~\ref{sec:gw150914} and the scaling run of Sec.~\ref{sec:scaling};
they are mapped to sections in that repository's \texttt{README.md}.  The
other 14 are 11 exploratory scans kept for provenance (not quoted here)
and 3 CarpetX smoke tests.  One supporting check quoted in
Sec.~\ref{sec:gw150914} --- the seven-level static-box functional check
--- was run from a production script that is not checked in.  The thorn drives both mesh drivers in the
Einstein Toolkit, \texttt{Carpet} (used for every result above) and
\texttt{CarpetX}, for which it provides explicit \texttt{OPTIONS: GLOBAL}
schedule entries because \texttt{CarpetX::decode\_mode} defaults to local
mode and would otherwise call the particle loop once per tile.

\section*{Statement on AI assistance}

During the preparation of this paper the author used AI assistants
(large language models) for debugging the code, drafting and revising
the text, and producing the figures.  All simulations, data and results
are the author's own; every number in the paper was checked by the
author, who takes full responsibility for its content.

\bibliographystyle{plainnat}
\bibliography{references}

\end{document}
